\section{Related Work}
\label{sec:related-work}

\subsection{From automated repair to repository-level agents}

Automated program repair (APR) established test-guided patch generation as a research area well before repository-level language-model agents \citep{monperrus2018bibliography,gazzola2019survey}. Defects4J supplied reproducible real faults for controlled Java studies \citep{just2014defects4j}; subsequent cross-benchmark experiments showed that apparent repairability can vary substantially with the benchmark and execution framework \citep{durieux2019repairthemall,ye2019quixbugs}. This history matters here because a repair statistic is conditional on both the task surface and the machinery that produces and validates candidates.

SWE-bench made executable issue resolution in real repositories a standard target for language-model evaluation \citep{jimenez2024swebench}. SWE-agent, Agentless, AutoCodeRover, and OpenHands illustrate how interfaces, localization procedures, sandboxes, and orchestration shape the path from issue text to candidate patch \citep{yang2024sweagent,xia2024agentless,zhang2024autocoderover,wang2025openhands}. Later benchmarks broadened the surface to live issues, multiple languages, long-horizon upgrades, and paid freelance tasks \citep{zhang2025swebenchlive,rashid2025swepolybench,thai2025sweevo,xu2026roadmapbench,miserendino2025swelancer}. SWE-smith further makes task construction and executable-environment generation explicit parts of the software-agent data pipeline \citep{yang2025swesmith}. Together, these systems expose many observable events before the final test result and make instrumentation part of the evaluation condition.

Recent work uses trajectories to recover process information omitted by resolve rate. AgentLens classifies passing trajectories by process quality and identifies ``lucky pass'' patterns \citep{sahoo2026agentlens}. TraceProbe normalizes tool trajectories and detects structural patterns such as search loops and skipped verification \citep{shu2026traceprobe}. Cross-framework analysis shows that the association between a behavioral signal and resolution can change direction across agent configurations \citep{ma2026samesignal}. Process evaluation also depends on the level of analysis: action prediction, task uncertainty, and step-level causal contribution answer different questions \citep{he2026process}.

These studies analyze what agents do over time. Our audit starts one level lower, at the measurement fields produced by the execution pipeline. It asks whether a logged field satisfies its stated invariant, then follows that construct to the verifier and through the available aggregation units.

\subsection{From weak test oracles to endpoint adequacy}

Test-based repair has long separated a patch that passes the available suite from a patch that satisfies the intended behavior. Generate-and-validate search spaces can contain orders of magnitude more test-passing incorrect patches than correct ones \citep{long2016searchspaces}. Empirical studies likewise document test-suite overfitting and show that repair reliability depends on the strength and composition of the test suite \citep{qi2015plausibility,smith2015overfitting,yi2018testsuite}. These results establish the first link in the historical chain relevant to our endpoint: a finite test oracle licenses a bounded operational claim.

Patch-correctness research then developed ways to probe that boundary. DiffTGen generates tests that distinguish candidate behavior from a reference implementation \citep{xin2017difftgen}. Behavior-similarity methods use executions to filter likely incorrect patches \citep{xiong2018patchcorrectness}. Comparative studies show that correctness assessors differ in precision, recall, and project sensitivity \citep{wang2020apca,le2019reliability}. Large-scale random testing against developer patches extends this assessment across repair tools while retaining an explicit reference-oracle assumption \citep{ye2021patchassessment}. Metamorphic testing provides another route when useful input transformations and output relations can be specified \citep{segura2016metamorphic}. In LLM code generation, EvalPlus shows that expanded tests can reject solutions accepted by the original benchmark and alter model rankings \citep{liu2023evalplus}.

Repository-level work continues this endpoint audit. BSG-VA replays an agent's validation command on buggy, candidate, and gold states \citep{xu2026validation}. Studies of agent-generated tests ask whether those tests support issue resolution or mainly provide feedback \citep{chen2026agenttests}. PatchDiff compares benchmark-passing patches with developer patches and reveals behavioral divergence after acceptance \citep{wang2025patchdiff}. The broader verification literature describes executable rewards as proxies for user intent whose faithfulness and robustness can change with task and policy capability \citep{wang2026verification}.

Our endpoint controls occupy a narrower role. They test whether the frozen verifier separates known-good, unrepaired, and selected near-reference candidates on the audited tasks. Reference-derived mutants are local probes rather than substitutes for the full distribution of real faults, a boundary established in mutation-testing research \citep{just2014mutants}. The controls therefore calibrate the endpoint used in this analysis without converting verifier pass into semantic adequacy.

\subsection{Aggregation, protocol, measurement, and reproducibility}

Pass@$k$ makes the number of sampled code candidates visible \citep{chen2021codex}. Recent work shows that this estimator can be misapplied when unit tests are treated as independent rollouts. Reliability@$k$ instead aggregates repeated complete attempts \citep{jiang2026beyond}. Our repeated-opportunity measure applies the same retained opportunities to governed action and verifier pass, allowing matched-budget comparison of their accumulation.

Evaluation validity also depends on task selection and protocol. LiveCodeBench uses temporal segmentation to reduce contamination and benchmark saturation \citep{jain2025livecodebench}. HackDetect audits agent exploitation of exposed artifacts and invalid scoring paths \citep{shao2026protocol}. An agent-safety benchmark audit shows that metric choice, target behavior, and model panel can change the meaning of a quoted score \citep{wang2026safetyaudit}. Broader analyses trace validity loss through task construction, simulation, and judgment \citep{caban2026measurement}.

The claim index is grounded in a longer software-measurement tradition. Property-based measurement work warns that different concepts can share one metric label unless the measured property is explicitly defined \citep{briand1996measurement}. Empirical-software-engineering guidance likewise treats units, context, design, and validity as reportable parts of an empirical result \citep{kitchenham2002guidelines,wohlin2012experimentation,ralph2018construct,sjoberg2023construct,ralph2020empiricalstandards}. Repository-mining research adds a provenance requirement: mutable sources, retrieval procedures, and archived derivatives determine whether a result can be reproduced \citep{gonzalez2012reproducibility}. Our five-coordinate index specializes these principles to repository-level repair by binding a statistic to its construct, opportunity design, support unit, endpoint interpretation, and instrumentation provenance.

\begin{table}[H]
\centering
\caption{Adjacent evaluation boundaries. The present study connects field integrity, criterion relation, opportunity, support, and endpoint scope in one retained execution corpus.}
\label{tab:positioning}
\small
\begin{tabularx}{\linewidth}{@{}L{0.21\linewidth}L{0.22\linewidth}L{0.16\linewidth}Y@{}}
\toprule
Work line & Audited object & Primary unit & Central question \\
\midrule
Trajectory diagnostics \citep{sahoo2026agentlens,shu2026traceprobe} & Search, edit, tool, and validation sequence & Trajectory & What process structure accompanies a resolved or unresolved run? \\
Cross-framework behavior \citep{ma2026samesignal} & Behavioral feature and outcome association & Agent configuration & Does a signal retain its direction across frameworks? \\
Validation evidence \citep{xu2026validation} & Agent-executed validation command & Validation event & Does a positive check distinguish buggy, candidate, and gold states? \\
Post-pass correctness \citep{wang2025patchdiff} & Benchmark-passing patch & Patch & Does a plausible patch match the behavior intended by the developer patch? \\
Repeated-rollout metrics \citep{jiang2026beyond} & Fully passing rollout under repeated attempts & Task--rollout set & Is repeated success aggregated with the intended sampling unit? \\
This study & Logged field, governed action, and verifier endpoint & Row, cell, task, family & Which claim is licensed after construct, budget, support, and endpoint are fixed? \\
\bottomrule
\end{tabularx}
\end{table}


\Cref{tab:positioning} locates the present study among these lines. Its unit spans a nested evidence chain: raw field, harmonized action construct, verifier endpoint, repeated cell, task, and task family. The contribution lies in auditing those boundaries together and attaching each reported value to an explicit claim index.
